\begin{examplebox}
\textbf{\textcolor{CExample}{例 1（基础）}}

\textbf{\textcolor{CExample}{题目：}} 设 $f(x) = 3x - 2$，解不等式 $1 < f(x) < 7$。

\textbf{\textcolor{CExample}{解题步骤：}}
\begin{description}[style=nextline, leftmargin=2.8em, labelsep=0.5em, itemsep=0.45em]
    \item[\textbf{第一步（写出原始）：}] $1 < 3x - 2 < 7$。
    \par\textbf{理由：}
    将 $f(x)$ 代入不等式。

    \item[\textbf{第二步（乘积形式）：}] 转化为乘积形式 $(3x-2-7)(3x-2-1) < 0$。
    \par\textbf{理由：}
    令 $M=7, N=1$，则乘积形式为 $(f(x)-M)(f(x)-N) < 0$。
    \[
    (3x-9)(3x-3) < 0
    \]

    \item[\textbf{第三步（解不等式）：}] 化简并解二次不等式。
    \[
    9(x-3)(x-1) < 0 \quad \Rightarrow \quad (x-3)(x-1) < 0
    \]
    \par\textbf{理由：}
    除以正数 9 不等号方向不变。解得 $1 < x < 3$。

    \item[\textbf{第四步（回代验证）：}] 将解集代回原始不等式验证。
    \par\textbf{理由：}
    确保解的正确性。当 $1 < x < 3$ 时，$3x-2$ 确实介于 1 和 7 之间。
\end{description}

\textbf{\textcolor{CExample}{关键结论：}} \textbf{不等式 $1 < 3x-2 < 7$ 的解集为 $x \in (1, 3)$。}

\medskip
\textbf{\textcolor{CExample}{例 2（稍变形）}}

\textbf{\textcolor{CExample}{题目：}} 设 $f(x) = x^2 + 1$，已知 $2 < f(x) < 10$，求 $x$ 的取值范围。

\textbf{\textcolor{CExample}{解题步骤：}}
\begin{description}[style=nextline, leftmargin=2.8em, labelsep=0.5em, itemsep=0.45em]
    \item[\textbf{第一步（绝对值形式）：}] 将不等式转化为绝对值形式。
    \par\textbf{理由：}
    令 $M=10, N=2$，则中点 $C = \frac{10+2}{2} = 6$，半长 $R = \frac{10-2}{2} = 4$。绝对值形式为
    \[
    |f(x) - 6| < 4
    \]
    即 $|x^2 + 1 - 6| < 4$。

    \item[\textbf{第二步（化简绝对值）：}] 化简得 $|x^2 - 5| < 4$。
    \par\textbf{理由：}
    计算 $f(x)-6 = x^2 + 1 - 6 = x^2 - 5$。

    \item[\textbf{第三步（解绝对值不等式）：}] 解 $|x^2 - 5| < 4$。
    \[
    -4 < x^2 - 5 < 4 \quad \Rightarrow \quad 1 < x^2 < 9
    \]
    \par\textbf{理由：}
    绝对值不等式等价于连不等式。分别解 $x^2 > 1$ 和 $x^2 < 9$。

    \item[\textbf{第四步（求 $x$ 范围）：}] 由 $1 < x^2 < 9$，得 $x^2 < 9$ 且 $x^2 > 1$，所以 $x \in (-3, -1) \cup (1, 3)$。
    \par\textbf{理由：}
    $x^2 < 9$ 推出 $-3 < x < 3$，$x^2 > 1$ 推出 $x < -1$ 或 $x > 1$，取交集。
\end{description}

\textbf{\textcolor{CExample}{关键结论：}} \textbf{满足 $2 < x^2+1 < 10$ 的 $x$ 取值范围为 $(-3, -1) \cup (1, 3)$。}
\end{examplebox}
